% Source: preamble
\documentclass[11pt,a4paper]{article}
\usepackage[margin=1in]{geometry}
\usepackage[T1]{fontenc}
\usepackage{lmodern,microtype}
\usepackage{amsmath,amssymb,amsthm,mathtools}
\usepackage{xcolor}
\usepackage[colorlinks=true,hypertexnames=false,linkcolor=blue!45!black,citecolor=blue!45!black,urlcolor=blue!45!black]{hyperref}
\hypersetup{pdftitle={High-Dimensional Gaussian Approximations over Rectangles without Covariance Restrictions},pdfsubject={Bounded independent summands; rates optimal up to logarithmic factors},pdfauthor={Human Collaborators: CCKK (Victor Chernozhukov, Denis Chetverikov, Kengo Kato, Yuta Koike); AI Collaborators: GPT-6 Sol and GPT-6 Astra}}
\setlength{\parindent}{1.2em}
\setlength{\parskip}{0.12em}
\emergencystretch=2em
\allowdisplaybreaks[2]
\raggedbottom
\clubpenalty=10000
\widowpenalty=10000
\newtheorem{theorem}{Theorem}
\newtheorem{proposition}{Proposition}
\newtheorem{lemma}{Lemma}
\newtheorem{corollary}{Corollary}
\theoremstyle{remark}
\newtheorem{remark}{Remark}
\newcommand{\E}{\mathbb E}
\newcommand{\Pp}{\mathbb P}
\newcommand{\R}{\mathbb R}
\newcommand{\cR}{\mathcal R}
\newcommand{\cF}{\mathcal F}
\newcommand{\cP}{\mathcal P}
\newcommand{\cD}{\mathcal D}
\newcommand{\Cov}{\operatorname{Cov}}
\newcommand{\Var}{\operatorname{Var}}
\newcommand{\tr}{\operatorname{tr}}
\newcommand{\op}{\mathrm{op}}
\newcommand{\norm}[1]{\left\lVert#1\right\rVert}
\newcommand{\rh}{\rho_{\mathcal R}}
\newcommand{\ind}{\mathbf 1}
\newcommand{\dd}{\,d}
\newcommand{\Hh}{H_2}
\newcommand{\Ga}{\gamma_r}
\newcommand{\ph}{\varphi_r}
\newcommand{\surf}{\mathcal J}

\newcommand{\logp}{\log_+}

\hypersetup{pdftitle={Supplementary Material: High-Dimensional Gaussian Approximations over Rectangles without Covariance Restrictions}}
\renewcommand{\thesection}{S.\arabic{section}}
\renewcommand{\theequation}{S.\arabic{equation}}
\renewcommand{\thetheorem}{S.\arabic{theorem}}
\renewcommand{\theproposition}{S.\arabic{proposition}}
\renewcommand{\thelemma}{S.\arabic{lemma}}
\renewcommand{\thecorollary}{S.\arabic{corollary}}
\renewcommand{\theremark}{S.\arabic{remark}}
\makeatletter
\renewcommand{\l@section}{\@dottedtocline{1}{0em}{2.4em}}
\renewcommand{\l@subsection}{\@dottedtocline{2}{2.4em}{3.4em}}
\makeatother


% Source: setup

\begin{document}
\begin{center}
{\Large Supplementary Material}\\[.4em]
{\large High-Dimensional Gaussian Approximations\\
over Rectangles without Covariance Restrictions}\\[.75em]
{\normalsize\textbf{Human Collaborators: CCKK}}\\[.2em]
{\normalsize Victor Chernozhukov, Denis Chetverikov, Kengo Kato, Yuta Koike}\\[.5em]
{\normalsize\textbf{AI Collaborators:} GPT-6 Sol and GPT-6 Astra}\\[.6em]
{\small Proposed editorial split, 20 October 2026}
\end{center}
\noindent
This supplement contains lower-bound proofs, auxiliary reductions, and
further calculations. All proofs of the main upper bounds are contained
in Appendices A--E of the main document. Section S.2 recalls only the
notation and estimates needed here. Sections S.3--S.5 give the
zero-skewness tail calculations, lower-bound arguments, and coordinate
reductions. Sections S.6--S.7 preserve optional calculations on the open
logarithmic question. The separate investigations note develops the
subsequent transport and rank-reduction calculations.
\tableofcontents

\section{Notation and statements used in the proofs}\label{sec:results}
We repeat the standing assumptions and result statements needed below.
This register makes the supplement independently readable; its numbering
is local to this document. The main paper contains all proofs of its
upper bounds, including the analytic estimates and their optimization.
Let \(\mathcal R\) be the axis-aligned rectangles, and write
\[
 \rh(U,V)=\sup_{A\in\mathcal R}
 |\Pp(U\in A)-\Pp(V\in A)|,\qquad
 M_h=\sup_A|\E F_{A,h}(W)-\E F_{A,h}(Z)|,
\]
where the rational selector \(F_{A,h}\) is defined in
Section~\ref{sec:smoothing}. The latter supremum also allows
inconsistent signed-face threshold systems.
\begin{equation}
 \begin{gathered}
 W=\sum_{i=1}^n\xi_i,\qquad
 \xi_1,\ldots,\xi_n\text{ independent and centered},\qquad
 \norm{\xi_i}_\infty\le\delta\quad\text{a.s.},\\
 \Sigma=\Cov(W),\qquad \Sigma_{jj}=1\quad(j\le p),
 \qquad Z\sim N(0,\Sigma).
 \end{gathered}
 \label{eq:setup}
\end{equation}

The rows are independent. Put
\[
L=\log(ep),\qquad c=1-\max_{j\ne k}|\Sigma_{jk}|,
\qquad c=1\text{ if }p=1,\qquad
D_L=1+\log_+\frac1{\delta\sqrt L},
\]
where \(\log_+x=\max\{0,\log x\}\), and define
\begin{equation}
 \kappa=\begin{cases}\min\{L,c^{-1}\},&c>0,\\L,&c=0,\end{cases}
 \qquad E_\kappa=1+\log_+\frac{\sqrt\kappa}{\delta L},
 \qquad \log_+x=\max\{0,\log x\}.
 \label{eq:parameters}
\end{equation}
\begin{theorem}[General bound]\label{thm:optimized}
Under \eqref{eq:setup},
\begin{equation}
 \rh(W,Z)\le C\min\!\left\{1,
 \delta\left[L^{3/2}+L\sqrt{\kappa E_\kappa}\right]\right\}.
 \tag{R*}\label{eq:optimized}
\end{equation}
\end{theorem}

\begin{theorem}[Cancellation of third moments]\label{thm:zero-skew}
Under \eqref{eq:setup}, suppose that
\begin{equation}
 \sum_i\E\xi_i^{\otimes3}=0,
 \quad\text{that is,}\quad
 \sum_i\E[\xi_{ij}\xi_{ik}\xi_{i\ell}]=0
 \quad(j,k,\ell\le p).
 \label{eq:zero-third}
\end{equation}
Then
\begin{equation}
 \rh(W,Z)\le C\min\{1,\delta L\sqrt{\kappa E_\kappa}\}.
 \tag{ZS}\label{eq:zero-skew}
\end{equation}
\end{theorem}

\begin{corollary}[Bounded observations]\label{cor:observations}
Let \(X_1,\ldots,X_n\in\R^p\) be independent and centered, with
\(\max_{i,j}|X_{ij}|\le K_n\) almost surely. Define
\[
 W_n=\frac1{\sqrt n}\sum_iX_i,\qquad
 \sigma_j^2=\frac1n\sum_i\E X_{ij}^2,
 \qquad \underline\sigma=\min_j\sigma_j>0,
 \qquad Z_n\sim N(0,\Cov(W_n)).
\]
Put \(\delta_n=K_n/(\underline\sigma\sqrt n)\), and calculate
\(c,\kappa,E_\kappa\) from \eqref{eq:parameters} using the correlation
matrix of \(W_n\) and the envelope \(\delta_n\). With \(L=\log(ep)\),
\begin{equation}
 \rh(W_n,Z_n)\le C\min\!\left\{1,
 \frac{K_n}{\underline\sigma\sqrt n}
 \left[L^{3/2}+L\sqrt{\kappa E_\kappa}\right]\right\}.
 \label{eq:observations}
\end{equation}
If \(\sum_i\E X_i^{\otimes3}=0\), the term \(L^{3/2}\) can be
omitted. In covariance rank one, the scalar bound
\(C\min\{1,\delta_n\}\) applies in either case.
\end{corollary}

\begin{corollary}[Dimension range under zero skewness]\label{cor:zero-threshold}
In Corollary~\ref{cor:observations}, suppose that
\(K_n/\underline\sigma\le B\), \(c\ge c_0>0\), and
\(\sum_i\E X_i^{\otimes3}=0\), with \(B\ge1\) and \(c_0\) fixed.
Then
\begin{equation}
 \rh(W_n,Z_n)\le C_{B,c_0}\min\!\left\{1,
 \frac L{\sqrt n}\sqrt{1+\log_+\frac{\sqrt n}{L}}\right\}.
 \label{eq:zero-fixed}
\end{equation}
Consequently, \(\log p=o(\sqrt n)\) guarantees convergence to zero.
This dimension threshold is sharp for the worst-case error over the
stated class, even when \(B=1\).
\end{corollary}

\begin{corollary}[Two simpler bounds]\label{cor:forms}
Under \eqref{eq:setup},
\begin{equation}
 \rh(W,Z)\le C\min\{1,\delta L^{3/2}\sqrt{D_L}\}.
 \label{eq:unrestricted}
\end{equation}
If \(c>0\), then also
\begin{equation}
 \rh(W,Z)\le C\min\!\left\{1,
 \delta\left[L^{3/2}+L\sqrt{D_L/c}\right]\right\}.
 \label{eq:gap-refined}
\end{equation}
\end{corollary}


The optimizing scale used in the main paper is
\begin{equation}
h=A\delta\sqrt{LE_\kappa/\kappa},
\label{eq:optimal-scale}
\end{equation}
with a sufficiently large absolute constant \(A\).


% Source: recalled_inputs

\section{Recalled notation and upper-bound estimates}
\label{sup:recalled}\label{sec:smoothing}\label{sec:support}
The statements in this section are recalled for the calculations below.
Their proofs are in Appendices A--D of the main document; no proof of a
main upper bound is deferred to this supplement.
For a signed coordinate system with distinct normals
\(\nu_\alpha\in\{\pm e_j\}\), the selector is
\[
F_{A,h}(x)=\left(1+\sum_\alpha
e^{(\nu_\alpha\cdot x-a_\alpha)/h}\right)^{-1}.
\]
Factor \(\Sigma=BB^\top\) on its rank-\(r\) support. Then
\begin{equation}
 \xi_i=B\zeta_i^0,\qquad
 \sum_i\E\zeta_i^0(\zeta_i^0)^\top=I_r,\qquad Z\stackrel d=BG,
 \quad G\sim N(0,I_r),
 \label{eq:support}
\end{equation}
and the support normals \(u_\alpha=B^\top\nu_\alpha\) have unit length.
Write \(\|z\|_U=\max_\alpha|u_\alpha\cdot z|\), so
\(\|\zeta_i^0\|_U\le\delta\). Below \(F\) denotes the selector on
the support, with
\(\pi_0=F\),
\(\pi_\alpha=Fe^{(u_\alpha\cdot x-a_\alpha)/h}\), and
\(\theta=1-F\).
The endpoint estimate is
\begin{equation}
 \E|F_{A,h}(W)-\ind_A(W)|+
 \E|F_{A,h}(Z)-\ind_A(Z)|\le Ch\sqrt L\,\kappa,
 \qquad h\ge\delta.
 \label{eq:endpoint-min}
\end{equation}
\subsection{The interpolation}
\label{an:sec:path}

We can now define the comparison path.  Bernoulli thinning removes the
bounded rows while a Gaussian component supplies the lost covariance.
The choice of amplitude below also makes the derivative of the cubic
row coefficient vanish at the bounded endpoint; this will help control
the start of the path.

Use the independent support rows \(\zeta_i^0\) of
\eqref{eq:support}, so
\[
 \E\zeta_i^0=0,\qquad
 \sum_i\E\zeta_i^0(\zeta_i^0)^\top=I_r,
 \qquad \norm{\zeta_i^0}_U\le\delta.
\]
For \(0\le\lambda\le1\), define
\begin{equation}
 \begin{gathered}
 a=1+\lambda,\qquad q=\frac{1-\lambda}{(1+\lambda)^2},
 \qquad B_{i,\lambda}\sim\operatorname{Bernoulli}(q),\\
 T_i=aB_{i,\lambda}\zeta_i^0,\qquad V_\lambda=\sum_iT_i,
 \qquad U_\lambda=V_\lambda+\sqrt\lambda G,
 \end{gathered}
 \label{an:path}
\end{equation}
where the Bernoulli variables, the source rows, and \(G\) are
independent.  Only the law at each \(\lambda\) is used, so no
coupling of Bernoulli variables across \(\lambda\) is required.
The total covariance of \(U_\lambda\) is \(I_r\), since
\(qa^2=1-\lambda\).  For a rectangle \(A\), fix its selector
\(F=F_{A,h}\) on the support and put
\[
 P_A(\lambda)=\E F(U_\lambda),\qquad H_i=U_\lambda-T_i.
\]
Let \(\zeta_i,\widetilde\zeta_i\) be two independent fresh copies of
\(\zeta_i^0\), independent of the actual path array and \(G\).

\begin{lemma}[Exact fourth-order identity]\label{an:lem:identity}
For \(0<\lambda<1\),
\begin{equation}
 P_A'(\lambda)
 =-\frac\lambda3\sum_i\E D^3F(U_\lambda)
                    [\zeta_i,\zeta_i,\zeta_i]
       +\sum_i(\mathcal A_i+\mathcal B_i+\mathcal C_i),
 \label{an:fourth-identity}
\end{equation}
where
\begin{align}
 \mathcal A_i
 &=\int_0^1 k_\lambda(s)
       \E D^4F(H_i+sa\zeta_i)
                 [\zeta_i,\zeta_i,\zeta_i,\zeta_i]\dd s,
 \label{an:A}\\
 k_\lambda(s)
 &=\frac a6(1-s)^2\{(3-\lambda)s-2\lambda\},
 \label{an:kernel}\\
 \mathcal B_i
 &=\frac{1-\lambda}{2}\int_0^1(1-s)
       \E D^4F(H_i+sa\widetilde\zeta_i)
          [\zeta_i,\zeta_i,\widetilde\zeta_i,\widetilde\zeta_i]
       \dd s,
 \label{an:B}\\
 \mathcal C_i
 &=\frac{\lambda(1-\lambda)}{3a}\int_0^1
       \E D^4F(H_i+sa\widetilde\zeta_i)
          [\zeta_i,\zeta_i,\zeta_i,\widetilde\zeta_i]\dd s.
 \label{an:C}
\end{align}
The three displayed fourth-order terms account for the entire remainder.
\end{lemma}
\begin{proposition}[Two bounds for the path derivative]
\label{an:prop:derivative}
There are sufficiently large absolute constants \(A,d_0\) such that,
for \(h\ge A\delta\), uniformly over rectangles \(A\),
\begin{equation}
 |P_A'(\lambda)|
 \le C\sqrt L\left(\frac{\lambda\delta}{h^2}
                         +\frac{\delta^2}{h^3}\right),
 \qquad 0<\lambda<1.
 \label{an:early}
\end{equation}
For \(d_0\delta^2L\le\lambda<1\), one also has
\begin{equation}
 |P_A'(\lambda)|
 \le CL^{3/2}\left(\delta+\frac{\delta^2}{\lambda h}\right).
 \label{an:late}
\end{equation}
More specifically, for the fourth-order remainder
\(R_{4,A,h}(\lambda)=\sum_i(\mathcal A_i+\mathcal B_i+\mathcal C_i)\),
\begin{equation}
 |R_{4,A,h}(\lambda)|\le\frac{C\delta^2}{\lambda h}L^{3/2},
 \qquad
 \left|\sum_i\mathcal C_i\right|
       \le\frac{C\delta^2}{h}L^{3/2}.
 \label{an:remainder}
\end{equation}
\end{proposition}

The common-base cubic estimate used below is
\begin{equation}
\left|\frac\lambda3\sum_i\E D^3F(U_\lambda)
[\zeta_i,\zeta_i,\zeta_i]\right|\le C\delta L^{3/2},
\qquad \lambda\ge d_0\delta^2L.
\label{an:cubic-bound}
\end{equation}
Appendix B of the main document proves the path identities and these
estimates, including the restoration of the actual translated row bases.


% Source: zero_skewness_tails
\section{Zero-skewness dimension thresholds and tail calculations}\label{zs:section}

For bounded observations and a fixed positive correlation gap, the
upper bound gives convergence whenever \(\log p=o(\sqrt n)\).
Independent Rademacher coordinates show why this condition is exact:
at the level of a Gaussian maximum, the fourth-order tail correction
has size \((\log p)^2/n\). We prove both statements below.

\begin{proof}[Sufficiency in Corollary~\ref{cor:zero-threshold}]
Use the standardized summands
\(\xi_{ij}=X_{ij}/(\sigma_j\sqrt n)\) from
Corollary~\ref{cor:observations}. We may take \(\delta=B/\sqrt n\),
because \(K_n/\underline\sigma\le B\). Their sum has unit marginal
variances and the same correlation matrix as \(W_n\). The diagonal
map preserves both rectangle distance and the vanishing third tensor.
Since \(c\ge c_0>0\), we have \(\kappa\le c_0^{-1}\), and
Theorem~\ref{thm:zero-skew} gives
\begin{equation}
 \rh(W_n,Z_n)\le C_{B,c_0}\min\left\{1,
 \frac L{\sqrt n}
 \sqrt{1+\log_+\frac{\sqrt n}{L}}\right\}.
 \label{zs:fixed-gap}
\end{equation}
Constants inside the positive logarithm are absorbed by
\(C_{B,c_0}\). If \(L=o(\sqrt n)\), the right-hand side tends to zero:
\(x\sqrt{1+\log(1/x)}\to0\) as \(x\downarrow0\).
This is equivalent to \(\log p=o(\sqrt n)\).
\end{proof}

The next example prevents uniform convergence at the boundary
\(\log p\asymp\sqrt n\). It settles the dimensional threshold;
sharpness of the error rate in \eqref{zs:fixed-gap} is a separate
question.

\begin{proposition}[Rademacher boundary obstruction]
\label{zs:rademacher-boundary}
Let \(\varepsilon_{ij}\), \(i\le n,j\le p_n\), be independent
Rademacher variables, and put
\(W_{n,j}=n^{-1/2}\sum_i\varepsilon_{ij}\).
Suppose \(\log p_n\sim a\sqrt n\) for some \(a>0\), and let
\(x_n\) be determined by \(p_n\{1-\Phi(x_n)\}=1\). Then
\begin{equation}
 \Pp\!\left(\max_{j\le p_n}W_{n,j}\le x_n\right)
 -\Pp\!\left(\max_{j\le p_n}Z_j\le x_n\right)
 \longrightarrow e^{-\exp(-a^2/3)}-e^{-1}>0,
 \label{zs:boundary-limit}
\end{equation}
where the \(Z_j\) are independent standard normals. These arrays have
\(B=1\), \(c=1\), and rowwise central symmetry, hence satisfy
\eqref{eq:zero-third}.
\end{proposition}

\begin{proof}
Write \(S_n=\sum_{i=1}^n\varepsilon_{i1}\), and let \(k_n\) be the
smallest integer with the same parity as \(n\) and
\(k_n>\sqrt n\,x_n\). Set \(t_n=k_n/n\). Mills' ratio implies
\[
 x_n^2\sim2\log p_n\sim2a\sqrt n,
 \qquad t_n=\frac{x_n}{\sqrt n}+O(n^{-1}).
\]
In particular \(t_n\to0\), \(nt_n^2\to\infty\), and
\(nt_n^6\to0\). Stirling's formula, applied to the exact binomial
mass \(b_n(k)=\Pp(S_n=k)\), gives
\begin{equation}
 b_n(k_n)=\sqrt{\frac2{\pi n}}(1-t_n^2)^{-1/2}
              e^{-nI(t_n)}(1+o(1)),
 \qquad
 I(t)=\frac{(1+t)\log(1+t)+(1-t)\log(1-t)}2.
 \label{zs:binomial-mass}
\end{equation}
The relative error is \(o(1)\), since both binomial cell counts are
of order \(n\).

To recover the tail from this mass, we must account for the spacing
of the binomial lattice. Adjacent admissible masses satisfy
\[
 \frac{b_n(k_n+2j+2)}{b_n(k_n+2j)}
 =\frac{n-k_n-2j}{n+k_n+2j+2}.
\]
Choose \(M_n\to\infty\) sufficiently slowly that
\(M_nt_n^2\to0\) and \(M_n^2/(nt_n^2)\to0\).
Uniformly for \(0\le j\le M_n/t_n\), expansion of the logarithms
of these ratios yields
\[
 \log\frac{b_n(k_n+2j)}{b_n(k_n)}
 =-2t_nj+O(t_n^3j+j/n+j^2/n)
 =-2t_nj+o(1).
\]
The displayed error follows by summing
\(\log(1-t_n-2\ell/n)-\log(1+t_n+(2\ell+2)/n)\)
over \(0\le\ell<j\); throughout this range the arguments stay
uniformly away from zero. For every larger admissible \(j\), each
adjacent mass ratio is at most
\((1-t_n)/(1+t_n+2/n)\le e^{-2t_n}\).
The remainder of the sum, beyond \(M_n/t_n\), is consequently at
most \(b_n(k_n)e^{-2M_n+O(t_n)}/(1-e^{-2t_n})\).
The initial geometric sum and this negligible remainder give
\begin{equation}
 \Pp(S_n\ge k_n)
 =\frac{b_n(k_n)}{2t_n}(1+o(1))
 =\frac{e^{-nI(t_n)}}{x_n\sqrt{2\pi}}(1+o(1)).
 \label{zs:binomial-tail}
\end{equation}

Now
\[
 I(t)=\frac{t^2}2+\frac{t^4}{12}+O(t^6),
 \qquad
 1-\Phi(x_n)=\frac{e^{-x_n^2/2}}{x_n\sqrt{2\pi}}(1+o(1)).
\]
The parity rounding contributes only
\(nt_n^2-x_n^2=O(x_n/\sqrt n)+O(n^{-1})=o(1)\), and
\(nt_n^4-x_n^4/n=o(1)\). Therefore \eqref{zs:binomial-tail} implies
\begin{equation}
 \log\frac{\Pp(S_n/\sqrt n>x_n)}{1-\Phi(x_n)}
 =-nI(t_n)+\frac{x_n^2}2+o(1)
 =-\frac{x_n^4}{12n}+o(1)
 \longrightarrow-\frac{a^2}3.
 \label{zs:tail-ratio}
\end{equation}
Thus \(p_n\Pp(S_n/\sqrt n>x_n)\to e^{-a^2/3}\). Independence
across coordinates gives
\[
 \Pp(\max_{j\le p_n}W_{n,j}\le x_n)
 =\{1-\Pp(S_n/\sqrt n>x_n)\}^{p_n}
 \longrightarrow e^{-\exp(-a^2/3)},
\]
whereas the corresponding normal probability is
\((1-1/p_n)^{p_n}\to e^{-1}\). This proves the proposition.
\end{proof}

\begin{proof}[Necessity in Corollary~\ref{cor:zero-threshold}]
If \(\log p_n=o(\sqrt n)\) fails, there is a subsequence on which
\(\log p_n\ge b\sqrt n\) for some \(b>0\). Choose any \(0<a<b\)
and restrict the rectangle to the first
\(q_n=\lfloor e^{a\sqrt n}\rfloor\le p_n\) coordinates, leaving
the other coordinate intervals unrestricted. The Rademacher arrays
of Proposition~\ref{zs:rademacher-boundary} then give a positive limiting
discrepancy along that subsequence. They belong to every class with
envelope bound \(B\ge1\), correlation-gap lower bound
\(0<c_0\le1\), and the aggregate zero-third-tensor assumption.
Together with \eqref{zs:fixed-gap}, this proves that
\(\log p_n=o(\sqrt n)\) is exactly the uniform convergence threshold
for these fixed-\(B,c_0\) classes.
\end{proof}

\subsection{What remains between the upper and lower rates}

Proposition~4.1 of Chernozhukov--Chetverikov--Koike~\cite{CCK}
assumes independent identically distributed scalar entries with a
finite exponential moment, unit variance, zero third moment, and
nonzero fourth cumulant. If
\[
 (\log p)^2/n\longrightarrow0,
 \qquad (\log p)^3/n\longrightarrow\infty,
\]
it states a positive limit superior for the maximum-distribution
error multiplied by \(n/(\log p)^2\). This gives \(L^2/n\) as a
lower-bound scale in this range. Rademacher entries meet these
assumptions, so the obstruction applies directly to the bounded
symmetric class considered here. The upper bound \eqref{zs:fixed-gap}
exceeds \(L^2/n\) by the factor
\[
 \frac{\sqrt n}{L}
 \sqrt{1+\log\frac{\sqrt n}{L}}
\]
when \(n^{1/3}\ll L\ll\sqrt n\). Consequently the zero-skewness
theorem settles the dimensional convergence threshold, while an
error-rate gap remains.

There is also a lattice obstruction in fixed dimension: a normalized
Rademacher sum has a central atom of order \(n^{-1/2}\), and the
Gaussian probability of an interval shrinking to that atom tends
to zero. Hence its rectangle distance is at least a constant times
\(n^{-1/2}\). A uniform unsmoothed \(L^2/n\) upper bound would
contradict this example. Theorem~4.2 of~\cite{CCK} obtains the faster
zero-skewness rate for a random vector with an independent
nondegenerate Gaussian component, under a rowwise mixed-third-moment
condition. That additional Gaussian smoothing is absent from the
bounded-row setting here. Its rate and optimality conclusions therefore
require assumptions beyond those imposed here.

The aggregate cancellation condition used in
Theorem~\ref{thm:zero-skew} is weaker than rowwise zero third tensors.
The fourth-order estimates used here leave room for improvement.
An optimal unsmoothed error bound must account for both the
fourth-cumulant correction and the lattice obstruction. Its precise
form remains open.


% Source: lower_bounds
\section{Lower bounds and optimality}
\label{opt:section}

Without a condition on third moments, the scale
\(\delta L^{3/2}\) is unavoidable even for independent coordinates. A small skewness correction in each coordinate is amplified
at the level of the maximum, where the relevant normal threshold is
about \(\sqrt{2\log p}\). We first make this argument precise, then
compare it with the upper bound to identify matching rates and sharp
convergence thresholds.

The lower bounds concern the worst error over a class of distributions;
a particular distribution can have a smaller error. Throughout this
section, \(\Phi\) and \(\varphi\) denote the standard normal distribution
function and density.

\subsection{A lower bound from skewness}
\label{opt:published}

In our notation, Proposition~2.1 of \cite{CCK} assumes \(B_n\ge2\) and
all three conditions
\begin{equation}
 r_n:=\frac{B_n(\log p_n)^{3/2}}{\sqrt n}\longrightarrow0,
 \qquad \frac{B_n^4}{\sqrt{\log p_n}}\longrightarrow0,
 \qquad \frac1{p_nr_n}\longrightarrow0.
 \label{opt:cck-hypotheses}
\end{equation}
It constructs independent scalar entries, identical in distribution
across rows and columns, with mean zero, variance one, and envelope
\(B_n\), for which \(\liminf_n\rh(W_n,Z_n)/r_n>0\). The resulting
covariance is \(I_{p_n}\), so the correlation gap is \(c=1\).
The discrepancy already occurs on common-threshold lower orthants,
which are rectangles.

The next theorem removes the last two restrictions in
\eqref{opt:cck-hypotheses}. Only the vanishing-rate condition and
\(p_n\to\infty\) remain. The proof uses the same two-point construction,
but controls its scalar tails uniformly in the envelope and compares
the coordinate products exactly.

\begin{theorem}[Skewness lower bound]
\label{opt:skew-lower}
Suppose \(p_n\to\infty\), \(B_n\ge2\), and
\[r_n=B_n(\log p_n)^{3/2}/\sqrt n\longrightarrow0.\]
There are i.i.d. centered vectors \(X_{n,1},\ldots,X_{n,n}\) with
independent coordinates, \(\Cov(X_{n,i})=I_{p_n}\), and
\(\norm{X_{n,i}}_\infty\le B_n\) almost surely. Put
\[W_n=n^{-1/2}\sum_iX_{n,i},\qquad Z_n\sim N(0,I_{p_n}).\]
Then
\begin{equation}
 \liminf_{n\to\infty}\frac{\rh(W_n,Z_n)}{r_n}
 \ge\frac{\sqrt2}{4e}>0.
 \label{opt:lower-constant}
\end{equation}
The envelope \(B_n\) is attained by the construction.
\end{theorem}

The scalar estimate below isolates the effect of skewness on a tail.
Its strict inequality matters because the lower example is a lattice
law. In particular, the estimate must remain valid as the threshold
passes through an atom.

\begin{lemma}[Uniform scalar moderate-tail estimate]
\label{opt:moderate-tail}
Let \(S=\sum_{i=1}^n\eta_i\), where the real variables \(\eta_i\) are
independent and centered, \(|\eta_i|\le\delta\) almost surely, and
\(\sum_i\E\eta_i^2=1\). Write
\(\gamma=\sum_i\E\eta_i^3\). There are absolute constants \(a,C>0\)
such that, whenever \(x\ge0\) and \(\delta(1+x)\le a\),
\begin{equation}
 \left|\log\frac{\Pp(S>x)}{1-\Phi(x)}
       -\frac{\gamma x^3}{6}\right|
 \le C\{\delta(1+x)+\delta^2x^4\}.
 \label{opt:moderate-tail-formula}
\end{equation}
The tail is strict. The conclusion includes lattice distributions and
has constants independent of the lattice spacing or the position of
\(x\) relative to its atoms.
\end{lemma}

\begin{proof}
Put \(K(t)=\sum_i\log\E e^{t\eta_i}\). For each summand, its
exponential tilt at \(t\ge0\) has likelihood ratio between
\(e^{-2\delta t}\) and \(e^{2\delta t}\). If \(\Var_t\) denotes variance
under this tilt, the independent-copy formula for variance gives
\[
 e^{-4\delta t}\Var(\eta_i)
 \le\Var_t(\eta_i)\le e^{4\delta t}\Var(\eta_i).
\]
Consequently, for a fixed small absolute \(a_0>0\),
\(e^{-4a_0}\le K''(t)\le e^{4a_0}\) on
\(0\le t\le a_0/\delta\). The centered tilted variable is bounded by
\(2\delta\), and \(\Var_t(\eta_i)\le\delta^2\). Its fourth cumulant
therefore has absolute value at most
\(7\delta^2\Var_t(\eta_i)\). Summing gives
\(\sup_{0\le t\le a_0/\delta}|K^{(4)}(t)|\le C\delta^2\).
Also \(|\gamma|\le\delta\). Taylor's theorem now yields, uniformly in
this local range,
\begin{align}
 K(t)&=t^2/2+\gamma t^3/6+O(\delta^2t^4),\notag\\
 K'(t)&=t+\gamma t^2/2+O(\delta^2t^3),\notag\\
 K''(t)&=1+\gamma t+O(\delta^2t^2).
 \label{opt:cumulant-expansions}
\end{align}

We choose the tilt so that its mean equals \(x\). To see that this
choice lies in the range just controlled, put \(t_0=a_0/\delta\).
Then
\[
 K'(t_0)=\int_0^{t_0}K''(u)\,du
 \ge e^{-4a_0}a_0/\delta.
\]
Choose \(a<e^{-4a_0}a_0/2\). The condition in the lemma implies
\(x<K'(t_0)\). Since \(K'(0)=0\) and \(K''>0\), there is a unique
\(t\in[0,t_0)\) with \(K'(t)=x\), including \(t=0\) when \(x=0\).
The bounds on \(K''\) show \(t\asymp x\) for \(x>0\).
Inverting \eqref{opt:cumulant-expansions}, and decreasing \(a\) if
necessary, gives
\begin{equation}
 t=x-\gamma x^2/2+O(\delta^2x^3),\qquad
 K(t)-tx=-x^2/2+\gamma x^3/6+O(\delta^2x^4).
 \label{opt:tilt-root}
\end{equation}
Let \(v^2=K''(t)\) and \(s=tv\). Expanding the square root in
\eqref{opt:cumulant-expansions} gives
\(v=1+\gamma x/2+O(\delta^2x^2)\), and hence
\begin{equation}
 s=x+O(\delta^2x^3).
 \label{opt:tilted-scale}
\end{equation}

Under the joint tilt
\(d\Pp_t=e^{tS-K(t)}d\Pp\), the summands remain independent and
\(T=(S-x)/v\) has mean zero and variance one. Furthermore,
\[
 \sum_i\E_t|\eta_i-\E_t\eta_i|^3
 \le2\delta\sum_i\Var_t(\eta_i)=2\delta v^2.
\]
Scalar Berry--Esseen and the upper and lower bounds on \(v\) imply
\(\sup_y|\Pp_t(T\le y)-\Phi(y)|\le C\delta\).
To keep track of any atom at the threshold, use the identity
\[
 \E_t[e^{-sT}\ind_{\{T>0\}}]
 =\Pp_t(T>0)-s\int_0^\infty e^{-sy}\Pp_t(T>y)\,dy,
\]
valid for \(s\ge0\), with the second term interpreted as zero when
\(s=0\). Comparison with
the same identity for a standard normal variable shows that
\begin{equation}
 \E_t[e^{-sT}\ind_{\{T>0\}}]=J(s)+O(\delta),\qquad
 J(u)=\int_0^\infty e^{-uy}\varphi(y)\,dy.
 \label{opt:strict-tail}
\end{equation}
The same \(O(\delta)\) estimate therefore covers an atom at \(T=0\).
The integral gives \(J(u)\asymp(1+u)^{-1}\) and
\(|(\log J)'(u)|\le C/(1+u)\) for \(u\ge0\).
The error in \eqref{opt:strict-tail} is thus relative
\(O(\delta(1+s))\); choosing \(a\) small makes it less than \(1/2\).
Also \eqref{opt:tilted-scale} gives
\(\log J(s)-\log J(x)=O(\delta^2x^2)\).
Finally,
\[
 \Pp(S>x)=e^{K(t)-tx}\E_t[e^{-sT}\ind_{\{T>0\}}],\qquad
 1-\Phi(x)=e^{-x^2/2}J(x).
\]
Taking logarithms and using \eqref{opt:tilt-root} proves
\eqref{opt:moderate-tail-formula}; the smaller
\(O(\delta^2x^2)\) term is absorbed into its stated remainder.
\end{proof}

\begin{proof}[Proof of Theorem~\ref{opt:skew-lower}]
Use mutually independent coordinates and rows with the two-point law
\begin{equation}
 \Pp(X=B_n)=\frac1{B_n^2+1},\qquad
 \Pp(X=-B_n^{-1})=\frac{B_n^2}{B_n^2+1}.
 \label{opt:two-point-law}
\end{equation}
It has mean zero, variance one, and third moment
\(\mu_{3,n}=B_n-B_n^{-1}\ge3B_n/4\). Put
\[
 \delta_n=B_n/\sqrt n,\qquad
 x_n=\Phi^{-1}(e^{-1/p_n}),\qquad
 t_n=1-\Phi(x_n),\qquad q_n=\Pp(W_{n,1}>x_n).
\]
Then \(x_n\sim\sqrt{2\log p_n}\), \(p_nt_n\to1\), and
\(\delta_nx_n^3\to0\). Lemma~\ref{opt:moderate-tail} gives
\begin{equation}
 q_n=t_n\{1+s_n+o(r_n)\},\qquad
 s_n=\frac{\mu_{3,n}x_n^3}{6\sqrt n}.
 \label{opt:tail-linearization}
\end{equation}
Indeed its remainder divided by \(r_n\) is
\(O(x_n^{-2}+\delta_nx_n)=o(1)\), and exponentiating introduces only
\(O(r_n^2)=o(r_n)\). Moreover
\(\liminf s_n/r_n\ge\sqrt2/4\).
By exact independence, the discrepancy on the common-threshold orthant is
\[
 \Pp(\max_jW_{n,j}\le x_n)-\Pp(\max_jZ_{n,j}\le x_n)
 =(1-q_n)^{p_n}-(1-t_n)^{p_n}.
\]
The mean value theorem makes this
\(-p_nt_n(1-u_n)^{p_n-1}\{s_n+o(r_n)\}\), for \(u_n\) between
\(q_n\) and \(t_n\). Since \(q_n/t_n\to1\), the leading prefactor
converges to \(e^{-1}\). Thus, more precisely,
\begin{equation}
 \Pp(\max_jW_{n,j}\le x_n)-e^{-1}
 =-\frac{\mu_{3,n}x_n^3}{6e\sqrt n}(1+o(1)).
 \label{opt:exact-product-asymptotic}
\end{equation}
The choice of \(x_n\) makes \((1-t_n)^{p_n}=e^{-1}\) exactly.
This avoids an additive \(O(p_n^{-1})\) error and hence a restriction
on \(p_nr_n\). The uniform scalar estimate removes the separate growth
restriction on \(B_n\). These two observations give the enlarged regime
in the theorem, and \eqref{opt:lower-constant} follows.
\end{proof}

\begin{remark}[The minimal envelope]
\label{opt:envelope-one}
The same construction works for \(B_n\ge b_0>1\), with lower constant
\(\sqrt2(1-b_0^{-2})/(3e)\). At \(B_n=1\), a centered unit-variance
marginal bounded by one is necessarily Rademacher, so its third moment
vanishes. The theorem does not assert this skewness lower bound at the
minimal envelope. Even at this envelope, symmetric marginals need not
imply the mixed-tensor cancellation required by Theorem~\ref{thm:zero-skew}; see
the example at the start of Section~\ref{zs:section}.
\end{remark}

\subsection{Matching rates at a fixed correlation gap}
\label{opt:matching}

To compare the upper bound with the scale forced by skewness,
divide it by \(\delta L^{3/2}\). Before truncation at one, the ratio is
\[
 1+\sqrt{\kappa E_\kappa/L}.
\]
Thus the upper bound has baseline order whenever
\(\kappa E_\kappa=O(L)\). For \(cL\ge1\), this condition reads
\[
 1+\log_+\frac1{\delta L\sqrt c}=O(cL).
\]
These are sufficient conditions for the second term to be absorbed
by the first. For a fixed gap \(c\ge c_0>0\), it suffices that
\begin{equation}
 1+\log_+(1/\delta)\le K L
 \label{opt:dimension-absorbs-log}
\end{equation}
for a fixed \(K\).

To make the worst-case statements precise, fix \(c_0\in(0,1]\), and let
\(M_n(c_0;B,p)\) be the supremum of \(\rh(W,Z)\) over independent
centered observations \(X_1,\ldots,X_n\in\R^p\) satisfying
\begin{equation}
 \begin{gathered}
 \max_i\norm{X_i}_\infty\le B\quad\text{a.s.},\qquad
 W=n^{-1/2}\sum_iX_i,\qquad
 \Sigma=n^{-1}\sum_i\E X_iX_i^{\mathsf T},\qquad
 \Sigma_{jj}=1,\\
 1-\max_{j\ne k}|\Sigma_{jk}|\ge c_0,
 \qquad Z\sim N(0,\Sigma),
 \end{gathered}
 \label{opt:minimax-class}
\end{equation}
with the usual convention for the gap when \(p=1\). This class allows
nonidentically distributed rows. Its lower examples can nevertheless
be chosen identically distributed.

\begin{corollary}[Matching fixed-gap minimax rate]
\label{opt:fixed-gap-minimax}
Fix \(c_0\in(0,1]\). Suppose \(p_n\to\infty\), \(B_n\ge2\),
\(r_n=B_n(\log p_n)^{3/2}/\sqrt n\to0\), and, for a fixed \(K\),
\[
 1+\log_+(\sqrt n/B_n)\le K\log(ep_n).
\]
There are constants \(a_1>0\) absolute and
\(a_2=a_2(c_0,K)<\infty\) such that, eventually,
\begin{equation}
 a_1r_n\le M_n(c_0;B_n,p_n)\le a_2r_n.
 \label{opt:matching-minimax}
\end{equation}
In particular, for every fixed \(a>0\), if \(p_n\asymp n^a\) and
\(B_n\ge2\) with \(B_n(\log n)^{3/2}/\sqrt n\to0\), then
\begin{equation}
 M_n(c_0;B_n,p_n)\asymp_{c_0,a}
 \frac{B_n(\log n)^{3/2}}{\sqrt n}.
 \label{opt:polynomial-dimension}
\end{equation}
The envelope may grow subject only to this vanishing-rate condition.
\end{corollary}

\begin{proof}
Apply Theorem~\ref{thm:optimized} with \(\delta_n=B_n/\sqrt n\).
On the class \eqref{opt:minimax-class}, \(\kappa\le c_0^{-1}\) and
\(E_\kappa\le C(c_0)+\log_+(1/\delta_n)\). The stated condition
therefore implies \(\kappa E_\kappa\le C(c_0,K)L_n\), giving the upper
bound. The identity-covariance examples in
Theorem~\ref{opt:skew-lower} belong to every such class and give the
lower bound, since \(\log(ep_n)\sim\log p_n\). For polynomial dimension,
\(\log_+(\sqrt n/B_n)\le\tfrac12\log n\), whereas
\(\log(ep_n)\sim a\log n\), so the required logarithmic condition holds
for every permitted growing envelope.
\end{proof}

\begin{remark}[A varying lower bound on the gap]
\label{opt:varying-gap-class}
More generally, replace the fixed \(c_0\) in the class definition by
any sequence \(c_n\in(0,1]\). Set
\[
 \kappa_n^0=\min\{L_n,c_n^{-1}\},\qquad
 E_n^0=1+\log_+\frac{\sqrt{\kappa_n^0}}{\delta_nL_n}.
\]
Under \(p_n\to\infty\), \(B_n\ge2\), \(r_n\to0\), and
\(\kappa_n^0E_n^0\le K L_n\), one has
\(M_n(c_n;B_n,p_n)\asymp_K r_n\). For the upper bound,
\(\kappa E_\kappa\) is nondecreasing in \(\kappa\), and every covariance
in the class has \(\kappa\le\kappa_n^0\). For the lower bound, the
identity covariance belongs to every class. Thus the worst error over
the covariance class has this rate; the argument does not determine the error at each prescribed covariance.
\end{remark}

The logarithmic condition \eqref{opt:dimension-absorbs-log} is needed
above to compare rates, but the convergence threshold has a simpler
form. To identify it, we first examine the two-point example at the
critical scale \(r_n\asymp1\).

\begin{proposition}[Critical maximum distribution]
\label{opt:critical}
For the independent coordinates \eqref{opt:two-point-law}, suppose
\(p_n\to\infty\), \(B_n\ge2\),
\[
 r_n=\frac{B_n(\log p_n)^{3/2}}{\sqrt n}\longrightarrow r\in(0,\infty),
 \qquad 1-B_n^{-2}\longrightarrow\alpha\in[3/4,1].
\]
At \(x_n=\Phi^{-1}(e^{-1/p_n})\),
\begin{equation}
 \Pp(\max_jW_{n,j}\le x_n)
 \longrightarrow
 \exp\!\left\{-\exp\!\left(\frac{\sqrt2}{3}\alpha r\right)\right\},
 \qquad \Pp(\max_jZ_{n,j}\le x_n)=e^{-1}.
 \label{opt:critical-maxcdf}
\end{equation}
\end{proposition}

\begin{proof}
Put \(\delta_n=B_n/\sqrt n\), \(q_n=\Pp(W_{n,1}>x_n)\), and
\(t_n=1-\Phi(x_n)\). Here \(\delta_nx_n^3=O(1)\), while
\(\delta_n(1+x_n)+\delta_n^2x_n^4\to0\).
Lemma~\ref{opt:moderate-tail} gives
\(q_n/t_n\to\exp(\sqrt2\alpha r/3)\). Since \(p_nt_n\to1\), exact
independence gives
\((1-q_n)^{p_n}\to\exp\{-\exp(\sqrt2\alpha r/3)\}\).
\end{proof}

\begin{theorem}[Exact uniform consistency threshold at a fixed gap]
\label{opt:consistency}
For every fixed \(c_0\in(0,1]\), every integer dimension sequence
\(p_n\ge1\), and every envelope sequence \(B_n\ge2\), put
\(\delta_n=B_n/\sqrt n\) and \(L_n=\log(ep_n)\). Then
\begin{equation}
 M_n(c_0;B_n,p_n)\longrightarrow0
 \quad\Longleftrightarrow\quad
 \delta_nL_n^{3/2}\longrightarrow0.
 \label{opt:consistency-iff}
\end{equation}
For a uniformly bounded envelope \(B_n\ge2\), this is equivalently
\(\log p_n=o(n^{1/3})\).
\end{theorem}

\begin{proof}
For sufficiency, Theorem~\ref{thm:optimized} bounds the error by a
constant depending only on \(c_0\) times
\[
 \delta_nL_n^{3/2}
 +t_n\sqrt{1+\log_+(C(c_0)/t_n)},\qquad t_n=\delta_nL_n.
\]
If the first term tends to zero, then \(t_n\to0\), so the second term
also tends to zero. No condition such as
\eqref{opt:dimension-absorbs-log} is needed for consistency.

For necessity, first suppose \(\delta_n\) does not tend to zero.
The binomial-atom argument in
Proposition~\ref{opt:rank-one-minimax} below, used in one coordinate
and independently in the others, gives an identity-covariance example
with error at least \(a_0\min\{1,\delta_n\}\). Thus the worst-case error
cannot tend to zero. Otherwise \(\delta_n\to0\). If
\(\delta_nL_n^{3/2}\) fails to tend to zero, pass to a subsequence on
which it is at least \(\varepsilon>0\). Then \(p_n\to\infty\), and
\[
 \log p_n\ge\tfrac12\varepsilon^{2/3}\delta_n^{-2/3}
\]
eventually. Set \(a=\varepsilon^{2/3}/4\) and
\(m_n=\lfloor\exp(a\delta_n^{-2/3})\rfloor\le p_n\).
Use \eqref{opt:two-point-law} independently in all \(p_n\) coordinates,
and consider a rectangle constraining only the first \(m_n\).
Pass, if needed, to a further subsequence on which
\(1-B_n^{-2}\to\alpha\in[3/4,1]\). For these \(m_n\) coordinates,
\(\delta_n(\log m_n)^{3/2}\to a^{3/2}>0\), so
Proposition~\ref{opt:critical} gives a strictly positive limiting
orthant discrepancy. The unconstrained coordinates do not change it.
This contradicts convergence of the worst-case error to zero and proves
necessity.
\end{proof}

\subsection{Rank-one optimality and singular embeddings}
\label{opt:rank-and-singularity}

\begin{proposition}[Nonasymptotic rank-one minimax rate]
\label{opt:rank-one-minimax}
For every \(n\ge1\), \(B\ge1\), and dimension \(p\ge1\), let
\(M_n^{(1)}(B,p)\) be the supremum of rectangle error over independent
centered observations with envelope \(B\), unit average marginal
variances, and covariance rank one. There are absolute constants
\(a_0,A_0>0\) such that
\begin{equation}
 a_0\min\{1,B/\sqrt n\}
 \le M_n^{(1)}(B,p)
 \le A_0\min\{1,B/\sqrt n\}.
 \label{opt:rank-one-minimax-rate}
\end{equation}
\end{proposition}

\begin{proof}
The upper bound is Proposition~\ref{red:rank-one}. For the lower bound,
let \(X_1,\ldots,X_n\) be independent with law
\eqref{opt:two-point-law}, replacing \(B_n\) by \(B\), and repeat the
same scalar in each coordinate, with any fixed choice of signs.
The normalized covariance has rank one and unit diagonal. The scalar
sum is an affine function of \(N\sim\operatorname{Bin}(n,\pi)\), where
\(\pi=(B^2+1)^{-1}\). Its binomial variance satisfies
\(v=n\pi(1-\pi)\le n/B^2\). Chebyshev's inequality gives probability
at least \(3/4\) to an interval of radius \(2\sqrt v\), which contains
at most \(4\sqrt v+2\) integers. Hence
\[
 \max_k\Pp(N=k)\ge\frac{3/4}{4\sqrt v+2}
 \ge a_0\min\{1,B/\sqrt n\}.
\]
Intervals shrinking around an atom of this size have Gaussian
probability tending to zero. Constrain one observed coordinate to such
an interval and leave the others unrestricted. The rectangle distance
is therefore at least this atom size. The argument includes \(B=1\),
all parities of \(n\), and the nonvanishing-error range \(B\gtrsim\sqrt n\).
\end{proof}

\begin{proposition}[A singular lower example with a small gap]
\label{opt:singular-embedding}
The leading lower rate in Theorem~\ref{opt:skew-lower} persists for
singular covariance with no equal or opposite coordinate pairs, and
with arbitrarily small positive correlation gap. The lower rate has no
additional factor depending on the gap.
\end{proposition}

\begin{proof}
Start with \(d\ge2\) independent coordinates of the construction in
Theorem~\ref{opt:skew-lower}. For any \(0<\epsilon\le1\), append the
coordinate
\[
 Y_i=\left(X_{i1},\ldots,X_{id},
       \frac{X_{i1}+\epsilon X_{i2}}{\sqrt{1+\epsilon^2}}\right).
\]
The covariance has rank \(d\) in dimension \(d+1\) and unit diagonal.
Its only nonzero off-diagonal correlations are
\(1/\sqrt{1+\epsilon^2}\) and
\(\epsilon/\sqrt{1+\epsilon^2}\), between the last coordinate and
the first and second, respectively. In particular no pair is equal or
opposite, and
\[
 c_\epsilon=1-\frac1{\sqrt{1+\epsilon^2}}>0,
 \qquad c_\epsilon\sim\epsilon^2/2\quad(\epsilon\downarrow0).
\]
The row envelope is
\(B_Y=B(1+\epsilon)/\sqrt{1+\epsilon^2}\in[B,\sqrt2B]\).
Rectangles ignoring the last coordinate reproduce every rectangle in
the original \(d\) coordinates, so
\(\rh(W_Y,Z_Y)\ge\rh(W_X,Z_X)\). Thus, along the vanishing-rate
sequences of Theorem~\ref{opt:skew-lower}, the lower rate is still an
absolute constant times
\((B_Y/\sqrt n)\{\log(d+1)\}^{3/2}\). One may choose
\(\epsilon=\epsilon_n\downarrow0\) arbitrarily fast. No coordinates
are removed by the reduction in Section~\ref{sec:coordinate-reduction}.
\end{proof}

\subsection{The remaining gaps}
\label{opt:remaining}

The matching regimes now have sharp dependence on sample size,
dimension, and envelope: Theorem~\ref{opt:skew-lower} supplies the
lower bound, and Corollary~\ref{opt:fixed-gap-minimax} supplies the upper
bound. Theorem~\ref{opt:consistency} identifies the fixed-gap convergence
threshold even outside those regimes. Rank one is also settled, for
all \(n,p,B\), by Proposition~\ref{opt:rank-one-minimax}.

The unrestricted factor \(\sqrt{D_L}\) and the additional dependence
on the correlation gap remain unresolved. The singular example above
forces \(\delta L^{3/2}\), but gives no factor involving
\(c^{-1/2}\), \(\sqrt{E_\kappa}\), or \(\sqrt{D_L}\). The
identity-covariance example likewise leaves open the optimal rate at a
prescribed covariance matrix or a prescribed nontrivial rank.
Proposition~\ref{red:coordinate-reduction} allows repeated or opposite
coordinates to be removed without changing rectangle distance.
This expresses the lower logarithm in the number of distinct coordinate
directions; it does not justify replacing that number by covariance
rank.

There is also an error-rate gap when the dimension grows too slowly
to satisfy \eqref{opt:dimension-absorbs-log}. The issue is uniformity
as dimension and covariance vary. For fixed \(p\) and fixed positive
definite \(\Sigma\), whitening and the classical convex-set
Berry--Esseen inequality give
\(\rh(W,Z)\le C_{p,\Sigma}\delta\). Indeed
\[
 \sum_i\E\norm{\Sigma^{-1/2}\xi_i}_2^3
 \le \sqrt p\,\norm{\Sigma^{-1/2}}_{\op}\delta
       \sum_i\E\norm{\Sigma^{-1/2}\xi_i}_2^2
 =p^{3/2}\norm{\Sigma^{-1/2}}_{\op}\delta.
\]
The constant here may depend on both \(p\) and the smallest
eigenvalue of \(\Sigma\), so it does not resolve the uniform question.

Finally, the discussion of near-optimal rates for unrestricted covariance
in the introduction of \cite{CCK} is a conjecture. Lower bounds for
strong sup-norm coupling, including the empirical-process comparison
mentioned there, do not by themselves give lower bounds for
rectangle-probability distance.


% Source: coordinate_reductions
\section{Exact coordinate reduction and the rank-one case}
\label{sec:coordinate-reduction}

A covariance matrix may have a zero correlation gap simply because two
coordinates are copies, with the same or opposite signs. Removing those
copies reduces the dimension in the bounds and leaves rectangle distance
unchanged. We prove this assertion first, then record the reduced bounds
and the scalar exception at rank one.

\begin{proposition}[Exact reduction of equal or opposite coordinates]
\label{red:coordinate-reduction}
Under \eqref{eq:setup}, define a relation on \(\{1,\ldots,p\}\) by
\[
 j\sim k\quad\Longleftrightarrow\quad |\Sigma_{jk}|=1.
\]
This is an equivalence relation.  Let \(C_1,\ldots,C_q\) be its classes,
choose \(j_a\in C_a\), and put
\[
 \begin{gathered}
 \varepsilon_j=\Sigma_{j,j_a}\in\{-1,1\}\quad(j\in C_a),
 \qquad \bar\Sigma=\Cov(\bar W),\\
 \bar\xi_i=(\xi_{i,j_a})_{a=1}^q,\quad
 \bar W=(W_{j_a})_{a=1}^q,\quad
 \bar Z=(Z_{j_a})_{a=1}^q.
 \end{gathered}
\]
Then, simultaneously for all applicable indices,
\begin{equation}
 \xi_{ij}=\varepsilon_j\xi_{i,j_a},\qquad
 W_j=\varepsilon_j W_{j_a},\qquad
 Z_j=\varepsilon_j Z_{j_a}
 \quad\text{a.s.}\quad(j\in C_a).
 \label{red:relations}
\end{equation}
The reduced summands satisfy the same normalized assumptions with the
same \(\delta\), and
\begin{equation}
 \operatorname{rank}(\bar\Sigma)=\operatorname{rank}(\Sigma),
 \qquad
 \rh(W,Z)=\rh(\bar W,\bar Z).
 \label{red:rank-distance}
\end{equation}
Consequently the relevant dimension and correlation gap after reduction
are
\begin{equation}
 L_* =\log(eq),\qquad
 c_*=
 \begin{cases}
  1-\displaystyle\max_{a\ne b}|\bar\Sigma_{ab}|,&q>1,\\[0.4em]
  1,&q=1,
 \end{cases}
 \qquad c_*>0.
 \label{red:parameters}
\end{equation}
Every rectangle bound proved under the normalized assumptions can
therefore be applied in dimension \(q\), with \(L_*\) and \(c_*\), and
then transferred to \((W,Z)\) without any loss.
\end{proposition}

\begin{proof}
For any \(j,k\) with \(|\Sigma_{jk}|=1\), write
\(\varepsilon=\Sigma_{jk}\).  Independence of the centered summands gives
\[
 0=\Var(W_j-\varepsilon W_k)
  =\sum_{i=1}^n\E(\xi_{ij}-\varepsilon\xi_{ik})^2.
\]
Every term is nonnegative, so
\(\xi_{ij}=\varepsilon\xi_{ik}\) almost surely for each \(i\).
In particular \(W_j=\varepsilon W_k\) almost surely.  Reflexivity and
symmetry of \(\sim\) are immediate.  If \(j\sim k\) and \(k\sim\ell\),
these identities imply
\(W_j=(\Sigma_{jk}\Sigma_{k\ell})W_\ell\), whence
\(|\Sigma_{j\ell}|=1\).  Thus \(\sim\) is transitive.  The same
zero-variance argument applied to \(Z_j-\varepsilon Z_k\) proves the
Gaussian identity in \eqref{red:relations}.  There are only finitely many
indices, so all the identities hold on common probability-one events.

The vectors \(\bar\xi_i\) retain independence, centering and the bound
\(\norm{\bar\xi_i}_\infty\le\delta\).  Also
\(\bar\Sigma_{aa}=1\).  Let \(T\) be the \(p\times q\) matrix defined by
\((Tx)_j=\varepsilon_j x_a\) for \(j\in C_a\), and let \(P\) select
the representative coordinates, so that \(PT=I_q\).  Then
\[
 \Sigma=T\bar\Sigma T^{\mathsf T},\qquad
 \bar\Sigma=P\Sigma P^{\mathsf T}.
\]
The two elementary rank inequalities prove rank equality.  Distinct
representatives belong to different equivalence classes, so
\(|\bar\Sigma_{ab}|<1\) for \(a\ne b\); this proves \(c_*>0\).

For completeness, the rectangle identity is exact even with unbounded
intervals and arbitrary choices of included endpoints.  Write
\(A=\prod_{j=1}^p I_j\), where each \(I_j\) is an interval, and define
\[
 B_A=\prod_{a=1}^q
       \left(\bigcap_{j\in C_a}\varepsilon_j I_j\right),
 \qquad
 \varepsilon I=\{\varepsilon x:x\in I\}.
\]
Each factor is an interval, possibly empty or a singleton; reflection
transports endpoint inclusion exactly and interchanges infinite
endpoints when the sign is negative.  Since \(T^{-1}A=B_A\),
\eqref{red:relations} yields
\[
 \Pp(W\in A)=\Pp(\bar W\in B_A),\qquad
 \Pp(Z\in A)=\Pp(\bar Z\in B_A).
\]
An empty factor makes both probabilities zero.  This proves
\(\rh(W,Z)\le\rh(\bar W,\bar Z)\).  Conversely, for any reduced rectangle
\(B=\prod_a J_a\), take \(I_{j_a}=J_a\) and \(I_j=\R\) at all other
coordinates.  Then \(T^{-1}A=B\), proving the reverse inequality.
If rectangles are initially defined using only finite closed intervals,
the same conclusion follows by monotone approximation: every interval
with any endpoint convention is an increasing union of finite closed
intervals (and empty intervals give zero discrepancy).  Continuity of
probability from below preserves the supremum in the definition of
rectangle distance.
\end{proof}

\begin{proposition}[Sharp scalar reduction at rank one]
\label{red:rank-one}
Under the assumptions of Proposition~\ref{red:coordinate-reduction},
\(\operatorname{rank}(\Sigma)=1\) if and only if \(q=1\).  In this case,
\begin{equation}
 \rh(W,Z)\le C\min\{1,\delta\}.
 \label{red:scalar-rate}
\end{equation}
Thus no dimension logarithm or smoothing logarithm is needed at rank
one, irrespective of the original number of coordinates.
\end{proposition}

\begin{proof}
If \(q=1\), rank preservation and \(\bar\Sigma=(1)\) give
\(\operatorname{rank}(\Sigma)=1\).  Conversely, a positive semidefinite
rank-one matrix with unit diagonal has the form
\(\Sigma=ss^{\mathsf T}\), where every \(s_j\in\{-1,1\}\).
Hence every pair of coordinates lies in the same equivalence class.

When \(q=1\), write \(S=\bar W=\sum_i\eta_i\), with
\(\eta_i=\xi_{i,j_1}\).  Then
\[
 \sum_i\E\eta_i^2=1,\qquad
 \sum_i\E|\eta_i|^3
 \le\delta\sum_i\E\eta_i^2=\delta.
\]
The scalar Berry--Esseen inequality gives
\[
 \sup_{t\in\R}|\Pp(S\le t)-\Phi(t)|\le C\delta,
\]
where \(\Phi\) is the standard normal distribution function.  The same
bound holds for strict inequalities by taking limits and using
continuity of \(\Phi\).  Subtracting the two endpoint distribution
functions bounds the discrepancy on every interval by \(2C\delta\),
including singleton intervals and arbitrary endpoint conventions.
Infinite endpoints are obtained by limits.  Now use
\eqref{red:rank-distance} and the trivial bound \(\rh\le1\).
\end{proof}

\begin{corollary}[Bounds in the reduced coordinates]\label{cor:reduced}
In the notation of Proposition~\ref{red:coordinate-reduction}, define
\[
 L_*=\log(eq),\qquad
 c_*=\begin{cases}
 1-\max_{a\ne b}|\bar\Sigma_{ab}|,&q>1,\\
 1,&q=1,
 \end{cases}
 \qquad \kappa_*=\min\{L_*,c_*^{-1}\}.
\]
Then \(c_*>0\), and
\begin{equation}
 \rh(W,Z)\le C\min\!\left\{1,
 \delta\left[L_*^{3/2}+L_*\sqrt{\kappa_*E_*}\right]\right\},
 \qquad E_*=1+\log_+\frac{\sqrt{\kappa_*}}{\delta L_*}.
 \label{eq:reduced-bound}
\end{equation}
The conclusions of Corollary~\ref{cor:forms} likewise hold with
\(L,c,D_L\) replaced by \(L_*,c_*,D_*\), where
\(D_*=1+\log_+(1/(\delta\sqrt{L_*}))\).
When \(q=1\), Proposition~\ref{red:rank-one} improves these bounds to
\(C\min\{1,\delta\}\).
Under \eqref{eq:zero-third}, the term \(L_*^{3/2}\) can be omitted
from \eqref{eq:reduced-bound} by Theorem~\ref{thm:zero-skew}.
\end{corollary}

\begin{proof}[Proof of Corollary~\ref{cor:reduced}]
Proposition~\ref{red:coordinate-reduction} gives independent centered
reduced rows bounded by \(\delta\), unit diagonal covariance, a positive
reduced correlation gap, and the exact distance equality
\eqref{red:rank-distance}. Apply Theorem~\ref{thm:optimized} and
Corollary~\ref{cor:forms} in dimension \(q\).
The rank-one improvement is Proposition~\ref{red:rank-one}.
If \eqref{eq:zero-third} holds, its reduced tensor is zero as well;
apply Theorem~\ref{thm:zero-skew} to omit the baseline term.
\end{proof}



% Source: further_calculations
\section{The signed remainder and the remaining logarithm}
\label{sup:status}

The remaining question is whether the logarithmic loss in
Theorem~\ref{thm:optimized} can be removed. The proposed bound is
\begin{equation}
 \rh(W,Z)\le C\min\{1,\delta L^{3/2}\}.
 \tag{T}\label{sup:target}
\end{equation}
The interpolation formula gives a precise way to state the missing
estimate. As we shall see, that estimate is equivalent to
\eqref{sup:target}, up to an error of the desired order; the reduction
alone supplies no further cancellation. Dimensions and normals may be
taken after the coordinate reduction in
Section~\ref{sec:coordinate-reduction}.

Choose sufficiently large absolute constants \(a_*,d_0\), and set
\[
 h=a_*\delta,\qquad \tau=d_0\delta^2L<\frac12.
\]
For a rectangle \(A\), let \(F\) be its selector on the covariance
support, and write \(P_A(\lambda)=\E F(U_\lambda)\). Keep the three
remainder terms from \eqref{an:A}, \eqref{an:B}, and \eqref{an:C}, and put
\begin{equation}
 R_{4,A,h}(\lambda)=\sum_i
       \{\mathcal A_i(\lambda)+\mathcal B_i(\lambda)+\mathcal C_i(\lambda)\},
 \qquad I_A=\int_\tau^1 R_{4,A,h}(\lambda)\dd\lambda.
 \label{sup:I}
\end{equation}
At this fixed smoothing scale, the missing estimate is
\begin{equation}
 \sup_{A\in\cR}|I_A|\le C\delta L^{3/2}.
 \tag{M}\label{sup:M}
\end{equation}
The absolute value is taken after integration in time. This order
matters: the pointwise absolute bound from the main proof loses a
logarithm.

\begin{proposition}[Reduction to the signed remainder]
\label{sup:reduction}
Uniformly over rectangles \(A\),
\begin{equation}
 I_A=\Pp(Z\in A)-\Pp(W\in A)+E_A,
 \qquad |E_A|\le C\delta L^{3/2}.
 \label{sup:exact-reduction}
\end{equation}
Consequently,
\begin{equation}
 \left|\sup_{A\in\cR}|I_A|-\rh(W,Z)\right|
 \le C\delta L^{3/2}.
 \label{sup:distance-reduction}
\end{equation}
Up to changes in its absolute constant, \eqref{sup:M} is equivalent to
\begin{equation}
 \sup_{A\in\cR}
 \left|\int_\tau^1\sum_i
       \{\mathcal A_i(\lambda)+\mathcal B_i(\lambda)\}\dd\lambda\right|
 \le C\delta L^{3/2}.
 \tag{$M_{AB}$}\label{sup:MAB}
\end{equation}
\end{proposition}

\begin{proof}
Integrating the exact identity \eqref{an:fourth-identity} gives
\begin{equation}
 I_A=P_A(1)-P_A(\tau)
       +\frac13\int_\tau^1\lambda
          \sum_i\E D^3F(U_\lambda)[\zeta_i,\zeta_i,\zeta_i]
          \dd\lambda.
 \label{sup:integrated-identity}
\end{equation}
The cubic estimate \eqref{an:cubic-bound} is
\[
 \left|\frac\lambda3
       \sum_i\E D^3F(U_\lambda)[\zeta_i,\zeta_i,\zeta_i]\right|
 \le C\delta L^{3/2},\qquad \lambda\ge\tau.
\]
Thus the last term of \eqref{sup:integrated-identity} has absolute
value at most \(C\delta L^{3/2}\). The early estimate
\eqref{an:early}, at \(h=a_*\delta\), gives
\[
 |P_A(\tau)-P_A(0)|
 \le C\sqrt L\left\{\frac{\delta\tau^2}{h^2}
                        +\frac{\delta^2\tau}{h^3}\right\}
 \le C\delta L^{3/2}(1+\delta^2L)
 \le C\delta L^{3/2}.
\]
The endpoint errors in \eqref{eq:endpoint-min} are also
\(O(\delta L^{3/2})\), since \(\kappa\le L\).
The path endpoints are \(BU_0\overset d=W\) and
\(BU_1\overset d=Z\), so these estimates prove
\eqref{sup:exact-reduction}. The inequality
\(\big||x|-|y|\big|\le|x-y|\), followed by a supremum, proves
\eqref{sup:distance-reduction}.

It remains to remove \(\mathcal C_i\). Its coefficient in
\eqref{an:C} contains an extra factor \(\lambda\), so the shifted-base estimate
that proves \eqref{an:remainder} yields
\begin{equation}
 \left|\sum_i\mathcal C_i(\lambda)\right|
 \le\frac{C\delta^2}{h}L^{3/2},\qquad \lambda\ge\tau.
 \label{sup:C-bound}
\end{equation}
At \(h=a_*\delta\), its time integral, even with an absolute value
at each time, is at most \(C\delta L^{3/2}\). Removing this controlled
term proves the equivalence with \eqref{sup:MAB}.
\end{proof}

When \(d_0\delta^2L\ge1/2\), the target follows from the trivial
bound \(\rh\le1\): indeed, \(L\ge1\) gives
\(\delta L^{3/2}\ge1/\sqrt{2d_0}\).
Thus proving \eqref{sup:M} in the stated range would prove the global
bound \eqref{sup:target}, and conversely. In rank one,
Proposition~\ref{red:rank-one} in
Section~\ref{sec:coordinate-reduction} already gives the stronger scalar
bound, so \eqref{sup:M} holds there.

What the present theorem gives is
\[
 \sup_A|I_A|
 \le C\delta\left[L^{3/2}+L\sqrt{\kappa E_\kappa}\right].
\]
This follows by combining Theorem~\ref{thm:optimized} with
Proposition~\ref{sup:reduction}. Improving it to \eqref{sup:MAB} requires
a new estimate for the signed integral.

\subsection{A time primitive and the derivative of the law}
\label{sup:primitive}

At \(\lambda=0\), the unresolved scalar kernels are
\[
 k_0(s)=\frac12s(1-s)^2,
 \qquad \frac12(1-s),\qquad 0\le s\le1.
\]
Both are nonnegative, so the scalar kernels alone do not cancel the
small-time singularity. Any such cancellation must also use the
averaged fourth derivatives.

An exact primitive for the \(\mathcal A_i\) kernel suggests an integration
by parts. It also shows why the resulting boundary term is not the
whole answer.
Put \(a_\lambda=1+\lambda\) and
\(q_\lambda=(1-\lambda)/a_\lambda^2\), and change variables
\(\vartheta=a_\lambda s\) in \eqref{an:A}. Then
\begin{equation}
 \mathcal A_i(\lambda)=\int_0^{a_\lambda}
  \partial_\lambda\left\{
    \frac{q_\lambda(a_\lambda-\vartheta)^3}{6}\right\}
  \E D^4F(H_{i,\lambda}+\vartheta\zeta_i)
       [\zeta_i,\zeta_i,\zeta_i,\zeta_i]\dd\vartheta.
 \label{sup:primitive-identity}
\end{equation}
The partial derivative holds \(\vartheta\) fixed. Indeed,
\[
 \left.\partial_\lambda
   \frac{q_\lambda(a_\lambda-\vartheta)^3}{6}
 \right|_{\vartheta=a_\lambda s}
 =\frac{k_\lambda(s)}{a_\lambda},
\]
where \(k_\lambda\) is \eqref{an:kernel}.

To carry out this integration by parts, define
\[
 g_i(\lambda,\vartheta)
 =\E D^4F(H_{i,\lambda}+\vartheta\zeta_i)
        [\zeta_i,\zeta_i,\zeta_i,\zeta_i],
 \qquad
 J_i(\lambda)=\int_0^{a_\lambda}
     \frac{q_\lambda(a_\lambda-\vartheta)^3}{6}
     g_i(\lambda,\vartheta)\dd\vartheta.
\]
The kernel vanishes at the moving upper endpoint, so
\begin{equation}
 \int_\tau^1\mathcal A_i(\lambda)\dd\lambda
 =J_i(1)-J_i(\tau)
 -\int_\tau^1\int_0^{a_\lambda}
   \frac{q_\lambda(a_\lambda-\vartheta)^3}{6}
   \partial_\lambda g_i(\lambda,\vartheta)
   \dd\vartheta\dd\lambda.
 \label{sup:primitive-correction}
\end{equation}
Here \(J_i(1)=0\), but the last term remains. It differentiates the law
of the deleted-row base \(H_{i,\lambda}\). If
\(M_i=\E\zeta_i\zeta_i^\top\), then
\(\Cov(H_{i,\lambda})=I_r-(1-\lambda)M_i\) and its covariance
derivative is \(M_i\). Differentiating the non-Gaussian rows also
produces mixed terms. Thus one cannot replace the integral of
\(\mathcal A_i\) by its boundary terms without controlling this change
of law. We have not bounded the additional signed contribution at the
scale required by \eqref{sup:MAB}.

\subsection{Where an early absolute value loses a logarithm}
\label{sup:conditional-example}

The next example isolates the effect of taking an absolute value
before averaging over the base variable. The total law will be
Gaussian at every time, yet the conditional absolute value will
produce a logarithm. Put
\[
 \beta(t)=\frac{e^t}{(1+e^t)^2},\qquad
 V_\lambda\sim N(0,1-\lambda),\qquad G\sim N(0,1),
\]
independently, and define
\[
 M_{\lambda,h}(v)=\E_G\left[(G^2-1)
       \beta\left(\frac{v+\sqrt\lambda G}{h}\right)\right].
\]
There exist absolute constants \(c,C,M>0\) such that, for all
sufficiently small \(h\),
\begin{equation}
 ch\log(1/h)
 \le\int_{Mh^2}^{1/2}
      \frac{\E|M_{\lambda,h}(V_\lambda)|}{\lambda}\dd\lambda
 \le Ch\log(1/h).
 \label{sup:conditional-log}
\end{equation}

To verify this, set \(\sigma=\sqrt\lambda\), let \(\varphi_1\)
be the standard normal density, and put
\(f(x)=(x^2-1)\varphi_1(x)\). A change of variables gives
\begin{equation}
 M_{\lambda,h}(v)
 =\frac h\sigma\int_\R
       f\left(\frac{hs-v}{\sigma}\right)\beta(s)\dd s.
 \label{sup:conditional-convolution}
\end{equation}
The function \(f\) is Lipschitz,
\(\int\beta=1\), and \(\int|s|\beta(s)\dd s<\infty\).
Uniformly for \(|v|\le\sigma/4\), the integral in
\eqref{sup:conditional-convolution} differs from \(f(-v/\sigma)\)
by at most \(Ch/\sigma\). The latter function is bounded above by
a negative absolute constant on this range. Choosing \(M\) large
therefore gives
\[
 M_{\lambda,h}(v)\le-\frac{ch}{\sigma},
 \qquad |v|\le\sigma/4,\quad \lambda\ge Mh^2.
\]
For \(\lambda\le1/2\),
\(\Pp(|V_\lambda|\le\sigma/4)\ge c\sigma\), proving
\(\E|M_{\lambda,h}(V_\lambda)|\ge ch\).
Conversely, the density of \(V_\lambda\) is bounded above by an
absolute constant for \(\lambda\le1/2\). Integrating the absolute
value in \eqref{sup:conditional-convolution} with respect to \(v\)
and using \(\int|f|<\infty\) gives
\(\E|M_{\lambda,h}(V_\lambda)|\le Ch\).
Integration of \(\dd\lambda/\lambda\) proves
\eqref{sup:conditional-log}.

Yet \(U=V_\lambda+\sqrt\lambda G\) is exactly standard Gaussian
for every \(\lambda\). Gaussian conditioning gives
\[
 \E[G^2-1\mid U]=\lambda(U^2-1),\qquad
 \E M_{\lambda,h}(V_\lambda)
 =\lambda\E[(U^2-1)\beta(U/h)].
\]
The expectation on the right is \(O(h)\), by the substitution
\(U=hs\) and integrability of \(\beta\). Hence
\begin{equation}
 \int_{Mh^2}^{1/2}
 \frac{|\E M_{\lambda,h}(V_\lambda)|}{\lambda}\dd\lambda
 \le Ch.
 \label{sup:signed-no-log}
\end{equation}
Thus the logarithm in \eqref{sup:conditional-log} disappears when
averaging precedes the absolute value. The example does not contradict
\eqref{sup:M} or \eqref{sup:target}; it identifies cancellation that a
bound on the conditional absolute value cannot retain.

\section{Two-face curvature in dimension two}
\label{sup:curvature}

A second approach keeps track of the angle between two faces. The
angle factor can be bounded by a positive curvature quantity, and in
dimension two we can bound its expectation uniformly. We give that
estimate here to record what the positive curvature argument supplies.
The later transport argument in the companion investigations~\cite{Cancellation}
controls the planar signed remainder; its extension to growing rank
remains open.

Let \(u_1,\ldots,u_m\in\R^r\) be unit vectors and let
\(a_1,\ldots,a_m\) be arbitrary thresholds. Define
\[
 \Psi(x)=h\log\left(1+\sum_{\alpha=1}^m
                    e^{(u_\alpha\cdot x-a_\alpha)/h}\right),
 \qquad H(x)=\nabla^2\Psi(x).
\]
Use the selector probabilities
\(\pi_0=F\),
\(\pi_\alpha=Fe^{(u_\alpha\cdot x-a_\alpha)/h}\), and set
\(u_0=0\). With
\[
 v=\sum_{\alpha=0}^m\pi_\alpha u_\alpha,\qquad
 C_\pi=\sum_{\alpha=0}^m\pi_\alpha
             (u_\alpha-v)(u_\alpha-v)^\top,
\]
we have \(C_\pi=hH\). For a symmetric matrix \(B\), put
\[
 \sigma_2(B)=\frac{(\tr B)^2-\tr(B^2)}2,
 \qquad
 \norm{z\wedge w}^2=\norm z^2\norm w^2-(z\cdot w)^2.
\]
The weighted-simplex identity is
\begin{equation}
 \sigma_2(C_\pi)
 =\sum_{0\le i<j<k\le m}
     \pi_i\pi_j\pi_k
       \norm{(u_j-u_i)\wedge(u_k-u_i)}^2.
 \label{sup:simplex}
\end{equation}
Indeed, each principal two-dimensional covariance determinant equals
the determinant of the augmented moment matrix of
\((1,u_{\alpha,s},u_{\alpha,t})\). The Cauchy--Binet formula expands
this determinant into weighted squared triangle determinants.
Summing over \(s<t\) gives \eqref{sup:simplex}. Retaining the
terms with \(i=0\) yields, in every support dimension,
\begin{equation}
 F\sum_{\alpha,\beta=1}^m
    \pi_\alpha\pi_\beta\{1-(u_\alpha\cdot u_\beta)^2\}
 \le 2\sigma_2(C_\pi)=2h^2\sigma_2(H).
 \label{sup:angle-curvature}
\end{equation}

\begin{proposition}[Rank-two curvature]
\label{sup:rank-two}
Let \(r=2\), and suppose
\[
 S=\sum_i\eta_i+G_0,\qquad \Cov(S)=I_2,
 \qquad \max_\alpha|u_\alpha\cdot\eta_i|\le\delta
 \quad\text{a.s.},
\]
where the \(\eta_i\) are independent and centered, and \(G_0\)
is an independent centered Gaussian vector. The Gaussian term is
optional and may be degenerate. If \(h\ge C\delta\), then,
uniformly over all thresholds and normal angles,
\begin{equation}
 \E\det H(S)\le C,
 \qquad
 \E\left[F(S)\sum_{\alpha,\beta=1}^m
   \pi_\alpha(S)\pi_\beta(S)
       \{1-(u_\alpha\cdot u_\beta)^2\}\right]\le Ch^2.
 \label{sup:rank-two-bound}
\end{equation}
\end{proposition}

\begin{proof}
We use a small-ball estimate to compare the expectation with a
Lebesgue integral, then bound that integral by the area of a gradient
image.

For the small-ball estimate, suppose
\(T=\sum_iZ_i+G\), where the centered rows are independent,
\(\norm{Z_i}_\infty\le\delta\), \(G\) is an independent centered
Gaussian vector, and \(V=\Cov(T)\) is positive definite. For every
square \(Q\) of side \(h\ge C\delta\),
\begin{equation}
 \Pp(T\in Q)\le\frac{Ch^2}{\sqrt{\det V}}.
 \label{sup:small-ball}
\end{equation}
To see this, use a product of squared sinc functions scaled by \(h\).
It gives a nonnegative integrable kernel bounded below on the square
centered at zero, with Fourier transform supported in
\(\{|t_j|\le C_0/h\}\) and bounded by \(C_1h^2\).
On this frequency box, \(|t\cdot Z_i|\) is smaller than an absolute
constant when \(h/\delta\) is sufficiently large. If \(Z_i'\)
is an independent copy, the local inequality
\(\cos x\le1-cx^2\) gives
\[
 |\varphi_{Z_i}(t)|^2
 =\E\cos\{t\cdot(Z_i-Z_i')\}
 \le1-c\,t^\top\Cov(Z_i)t.
\]
Independence and the Gaussian characteristic function imply
\( |\varphi_T(t)|\le e^{-c't^\top Vt}\) on this box.
Fourier inversion of the translated kernel therefore bounds the
probability by
\[
 Ch^2\int_{\R^2}e^{-c't^\top Vt}\dd t
 =\frac{Ch^2}{\sqrt{\det V}},
\]
proving \eqref{sup:small-ball}.

We now apply \eqref{sup:small-ball} in coordinates chosen from the
normals. If all the normals are collinear, \(\det H=0\), so both
conclusions are immediate. Otherwise choose \(u_a,u_b\) maximizing
\(d=|\det(u_a,u_b)|>0\). Cramer's rule and maximality show that
every normal has the representation
\[
 u_\alpha=A_\alpha u_a+B_\alpha u_b,
 \qquad |A_\alpha|,|B_\alpha|\le1.
\]
Tile the coordinates \((u_a\cdot x,u_b\cdot x)\) by half-open
\(h\)-squares, and denote their inverse images by \(Q_\ell\).
Each has area \(|Q_\ell|=h^2/d\). The transformed rows
\((u_a\cdot\eta_i,u_b\cdot\eta_i)\) have maximum norm at most
\(\delta\), and the covariance determinant of
\((u_a\cdot S,u_b\cdot S)\) is \(d^2\). Thus
\eqref{sup:small-ball} gives
\begin{equation}
 \Pp(S\in Q_\ell)\le C|Q_\ell|.
 \label{sup:cell-probability}
\end{equation}

If \(x,y\in Q_\ell\), then
\(|u_\alpha\cdot(x-y)|\le2h\) for every \(\alpha\).
All selector probabilities at \(x\) and \(y\) are consequently
within a fixed multiplicative factor. In two dimensions,
\eqref{sup:simplex} is a positive expansion of \(h^2\det H\).
The product of three probabilities in each term changes by only a
fixed factor, whence
\(\sup_{Q_\ell}\det H\le C\inf_{Q_\ell}\det H\).
Together with \eqref{sup:cell-probability}, this implies
\begin{equation}
 \E\det H(S)
 \le C\sum_\ell |Q_\ell|\inf_{Q_\ell}\det H
 \le C\int_{\R^2}\det H(x)\dd x.
 \label{sup:curvature-integral}
\end{equation}

It remains to bound the integral. The vectors \(0,u_a,u_b\) affinely
span \(\R^2\), and all selector probabilities are positive. Therefore
\(H\) is positive
definite everywhere. In particular \(\nabla\Psi\) is injective:
the integral of the Hessian along the line between two distinct points
has a strictly positive quadratic form in their difference.
The image of \(\nabla\Psi\) lies in
\(\operatorname{conv}\{0,u_1,\ldots,u_m\}\), a subset of the unit
disk. The area formula gives
\[
 \int_{\R^2}\det H(x)\dd x
 =|\nabla\Psi(\R^2)|\le\pi.
\]
This and \eqref{sup:curvature-integral} prove the first assertion.
The second follows from \eqref{sup:angle-curvature}, since
\(\sigma_2(H)=\det H\) in dimension two.
\end{proof}

Proposition~\ref{sup:rank-two} alone does not bound the signed
\(\mathcal A+\mathcal B\) remainder by a positive curvature quantity.
The later facewise transport argument in~\cite{Cancellation} supplies
a planar signed-remainder estimate through additional identities.
A uniform extension to growing rank remains open.
In higher rank, applying the grid argument to two-dimensional slices
would require small-ball control conditional on the complementary
coordinates. Independence of the original rows does not supply that
conditional estimate.


% Source: references
\begin{thebibliography}{9}
\bibitem{Cancellation}
\emph{Cancellation and rank reduction in rectangle CLTs:
transport identities, planar bounds, and rank-three tests}.
Companion research note, 20 October 2026.

\bibitem{AL}
P.~M. Aronow and P. Lopatto (2026).
\emph{Optimal central limit theorem for bounded random variables in high dimensions}.
Manuscript, arXiv:2609.07067v1.
\url{https://arxiv.org/abs/2609.07067v1}.

\bibitem{Bounded}
\emph{A rectangle CLT for bounded summands: correlation gaps and
unrestricted covariance}.
Preceding research note, October 19, 2026.
\texttt{Rectangle\_CLT\_Bounded\_Covariance.pdf}.

\bibitem{Integrated}
\emph{On the integrated remainder (M): a rank-free square-root
logarithmic bound}.
Preceding research note, October 19, 2026.
\texttt{Rectangle\_CLT\_Integrated\_Remainder.pdf}.

\bibitem{Gumbel}
\emph{The rational rectangle smoother and the Gumbel trick}.
Companion note.
\texttt{Rational\_Smoother\_and\_Gumbel\_Trick.pdf}.

\bibitem{Kane}
D. Kane (2014).
The average sensitivity of an intersection of half spaces.
\emph{Research in the Mathematical Sciences} \textbf{1}, Article 13.
\href{https://doi.org/10.1186/s40687-014-0013-6}{doi:10.1186/s40687-014-0013-6}.

\bibitem{CCK}
V. Chernozhukov, D. Chetverikov, and Y. Koike (2021).
\emph{Nearly optimal central limit theorem and bootstrap approximations
in high dimensions}.
arXiv:2012.09513v2, May 12, 2021.
\url{https://arxiv.org/abs/2012.09513v2}.

\bibitem{Local}
V. Chernozhukov, D. Chetverikov, K. Kato, and Y. Koike,
with AI collaborators (2026).
\emph{High-dimensional Gaussian approximation over rectangles under
local nondegeneracy}. Companion working manuscript supplied
20 October 2026.

\end{thebibliography}
\end{document}


